\subsection{Résultats}

\subsubsection{Temps de propagation}

Nous mesurons donc 2 signaux sur l'oscilloscope (visible sur \ref{fig:r0} : l'impulsion du generateur (CH1) et l'impulsion passant par le cable coaxial de 100m (CH2).

\begin{figure}[ht]
	\centering
	\includegraphics[scale=0.6]{electrique/r5.png}
	\label{fig:r0}
	\caption{Mesure des impulsions de largeur (width) $w$ = 210 $ns$}
\end{figure}

Le delai entre les deux signaux pour une resistance nulle R = 0 $\ohm$ $\delta$$t$ = $|$$t_1$ - $t_2$$|$ est : 

\[
        \mathrm{\delta t} = | 
\]

Que nous avons extrait sour \ref{fig:r0} 
Nous avons mesure grace au pont RLC les valeurs suivantes pour la capacite lineique $c*$ et $l*$ :

\[
	\mathrm{c*} = 10,1 nF \qquad
	\mathrm{l*} = 31,0 \mu H
\]

Ainsi comme nous connaissons la relation liant c* et l* avec l'impedance characteristique :
\[
	\mathrm{Z_c} = \sqrt{\frac{l*}{c*}} = \sqrt{\frac{31,0 \cdot 10^-6}{10,1 \cdot 10^-9}} = 55,5 \pm 0,3
\]

Pour les incertitudes :
\[
\Delta Z_c = \frac{1}{2} \cdot Z_c \sqrt{\left(\frac{\Delta L^*}{L^*}\right)^2 + \left(\frac{\Delta C^*}{C^*}\right)^2}
\]

Donc :
\[
\Delta Z_c = 55,5 \cdot \frac{1}{2} \sqrt{\left(\frac{0,1}{31,0}\right)^2 + \left(\frac{0,1}{10,1}\right)^2} \approx 0,3~\Omega
\]

subsubsection{Un deuxieme resultat}

